% ============================================================
% SECTION 4 — Results
% ============================================================

\section{Results}
\label{sec:results}

\subsection{Posterior constraints on the CCR parameters}

Table~\ref{tab:constraints} summarises the marginalised posterior
constraints on all model parameters. The six standard $\Lambda$CDM
parameters are recovered in excellent agreement with the Planck 2018
baseline results~\cite{Planck:2018params}, confirming that the CCR
modification does not introduce tension at high multipoles.

\begin{table}[t]
\centering
\begin{tabular}{lcc}
\toprule
Parameter & CCR constraint (68\% CL) & Planck 2018 $\Lambda$CDM \\
\midrule
$\log_{10}(\ln D)$ & $14.0 \pm 1.2$ & --- \\
$\alpha$ & $3.4 \pm 1.4$ & --- \\
$k_c\;[\mathrm{Mpc}^{-1}]$ & $0.00048^{+0.00007}_{-0.00051}$ & --- \\
$H_0\;[\mathrm{km\,s^{-1}\,Mpc^{-1}}]$ & $67.24 \pm 0.58$ & $67.36 \pm 0.54$ \\
$\Omega_b h^2$ & $0.02236 \pm 0.00014$ & $0.02237 \pm 0.00015$ \\
$\Omega_c h^2$ & $0.1202 \pm 0.0013$ & $0.1200 \pm 0.0012$ \\
$\tau$ & $0.054 \pm 0.005$ & $0.054 \pm 0.007$ \\
$10^9 A_s$ & $2.101 \pm 0.023$ & $2.100 \pm 0.030$ \\
$n_s$ & $0.9647 \pm 0.0042$ & $0.9649 \pm 0.0042$ \\
\bottomrule
\end{tabular}
\caption{Marginalised parameter constraints from the CCR MCMC analysis
  (68\% CL), compared with Planck 2018 $\Lambda$CDM baseline
  values~\cite{Planck:2018params}. The CCR parameters
  $\log_{10}(\ln D)$ and $\alpha$ have no $\Lambda$CDM counterpart.}
\label{tab:constraints}
\end{table}

The posterior on $\log_{10}(\ln D)$ is broad and approximately uniform
across the prior range $[12,\,16]$, indicating that the Planck data
do not strongly constrain the Hilbert space dimension
(figure~\ref{fig:corner_ccr}). The marginal posterior on $\alpha$ shows
a mild preference for larger values ($\alpha \gtrsim 3$), corresponding
to a relatively sharp suppression profile, though the constraint is
weak. The joint posterior in the $(\log_{10}(\ln D),\,\alpha)$ plane
exhibits a characteristic exclusion region at small $\ln D$ and small
$\alpha$, where the combination of a low cutoff scale with gradual
suppression would produce excessive power deficit incompatible with the
observed low-$\ell$ spectrum.

\begin{figure}[t]
\centering
\includegraphics[width=0.55\textwidth]{corner_ccr_params.png}
\caption{Joint posterior distribution in the
  $(\log_{10}(\ln D),\,\alpha)$ plane. Dark and light shading denote the
  68\% and 95\% credible regions, respectively. The exclusion of the
  lower-left corner (small $\ln D$, small~$\alpha$) reflects the
  data's rejection of strong, gradual suppression profiles.}
\label{fig:corner_ccr}
\end{figure}

\begin{figure}[t]
\centering
\includegraphics[width=\textwidth]{corner_all_params.png}
\caption{Full joint posterior for all sampled parameters. The CCR
  parameters $\log_{10}(\ln D)$ and $\alpha$ (top two rows) show no
  correlations with the standard $\Lambda$CDM parameters, confirming
  that the CCR modification does not shift or degrade the precision
  of cosmological constraints.}
\label{fig:corner_all}
\end{figure}

\subsection{Model comparison}
\label{sec:model_comparison}

The best-fit CCR model achieves $\chi^2_{\min} = 1003.4$, with best-fit
parameters $\log_{10}(\ln D) = 15.2$ and $\alpha = 1.8$. This best-fit
point lies in the $\Lambda$CDM-like regime
($\ln D = 1.6 \times 10^{15}$, corresponding to
$k_c = 4.6 \times 10^{-5}\;\mathrm{Mpc}^{-1}$), where the suppression
is negligible across all observable multipoles.

We compare the best-fit $\chi^2$ values in two distinct regions of
parameter space: the $\Lambda$CDM-like region
($\log_{10}(\ln D) > 15$) and the active-cutoff region
($\log_{10}(\ln D) < 14$, where
$k_c > 2 \times 10^{-4}\;\mathrm{Mpc}^{-1}$ produces measurable
suppression). The difference is
\begin{equation}
\Delta \chi^2 = \chi^2_{\rm cutoff} - \chi^2_{\Lambda{\rm CDM}} = 0.13\,,
\label{eq:deltachi2}
\end{equation}
indicating that the active cutoff region provides an essentially
identical fit to the $\Lambda$CDM-like region.

The Savage--Dickey density ratio yields
\begin{equation}
\ln B_{01} = 0.06\,,
\end{equation}
where $B_{01}$ is the Bayes factor for $\Lambda$CDM relative to CCR.
This value is effectively zero on the Jeffreys scale, indicating no
statistical preference for either model.

The information criteria yield
\begin{equation}
\Delta\mathrm{AIC} = \mathrm{AIC}_{\rm CCR} - \mathrm{AIC}_{\Lambda{\rm CDM}} = +4.0\,,
\qquad
\Delta\mathrm{BIC} = \mathrm{BIC}_{\rm CCR} - \mathrm{BIC}_{\Lambda{\rm CDM}} = +13.0\,.
\end{equation}
The positive $\Delta$AIC and $\Delta$BIC reflect the Occam penalty for
the two additional CCR parameters, which do not yield a compensating
improvement in fit quality. By the Jeffreys scale for
BIC~\cite{Kass:1995}, $|\Delta\mathrm{BIC}| > 10$ constitutes strong
evidence against the more complex model. This is unsurprising given the
broad posterior and null detection: our claim is not model preference
but model viability and forecasted detectability. The BIC penalty
reflects the expected Occam cost of two weakly constrained
parameters~\cite{Liddle:2007}, consistent with the data being
insufficiently informative to distinguish between the models at the
current sensitivity.

\subsection{Decoupling from standard cosmological parameters}
\label{subsec:decoupling}

A key result of this analysis is the absence of correlations between the
CCR parameters and the standard cosmological parameters. The joint
posterior (figure~\ref{fig:corner_all}) shows that $\log_{10}(\ln D)$
and $\alpha$ are effectively uncorrelated with $H_0$, $\Omega_b h^2$,
$\Omega_c h^2$, $\tau$, and~$n_s$. This decoupling is physically
expected: the CCR suppression affects only modes with
$k \lesssim k_c \sim 10^{-4}\;\mathrm{Mpc}^{-1}$, corresponding to
$\ell \lesssim 10$, while the standard parameters are primarily
constrained by the acoustic peak structure at $\ell \gtrsim 30$.

This decoupling has an important practical implication: the CCR
framework can be tested independently of the standard cosmological
parameters, and its inclusion does not shift or degrade the precision
of $\Lambda$CDM constraints.

\subsection{\texorpdfstring{$\Lambda$CDM}{ΛCDM} recovery}
\label{subsec:lcdm_recovery}

As a validation test, we verify that in the limit $\ln D \to \infty$
the CCR-modified $C_\ell$ spectrum recovers the standard $\Lambda$CDM
prediction. For $\log_{10}(\ln D) = 16$ the fractional difference
$|\Delta C_\ell / C_\ell^{\Lambda{\rm CDM}}|$ is below
$2.2 \times 10^{-4}$ at all multipoles $\ell = 2$--$2500$, confirming
that the CCR modification is a smooth, well-controlled deformation
of the standard model with a well-defined $\Lambda$CDM limit.

% ============================================================
% Section 4.5: Comparison with Planck Low-ℓ Data
% Insert after Section 4.4 (ΛCDM recovery)
% ============================================================

\subsection{Comparison with Planck low-\texorpdfstring{$\ell$}{ℓ} data}
\label{subsec:cl_comparison}

Figure~\ref{fig:cl_data_comparison} presents a direct comparison of the
Planck 2018 Commander low-$\ell$ TT power spectrum~\cite{Planck:2018likelihood}
with the CCR predictions for three representative points in parameter space:
(i)~the MCMC best-fit (log$_{10}(\ln D) = 15.2$, $\alpha = 1.8$), which lies
in the $\Lambda$CDM-like regime with negligible suppression;
(ii)~the fiducial CCR model ($\ln D = 2\times 10^{14}$, $\alpha = 2$),
which produces a modest $\sim 15\%$ suppression at $\ell = 2$; and
(iii)~an active-cutoff model ($\ln D = 2.2\times 10^{13}$, $\alpha = 2$),
corresponding to the Contaldi et al.\ best-fit cutoff scale
$k_c \simeq 4.9\times 10^{-4}\,\mathrm{Mpc}^{-1}$, which yields
$\sim 50\%$ suppression at the quadrupole.

All three models are visually consistent with the data, as the Planck
data points scatter broadly within the cosmic variance envelope. The
best-fit CCR model is indistinguishable from $\Lambda$CDM at all
multipoles, with a maximum fractional deviation of $8.6\times 10^{-3}$.

Figure~\ref{fig:cl_residuals} shows the fractional residuals
$(\mathcal{D}_\ell^{\rm CCR} - \mathcal{D}_\ell^{\Lambda\rm CDM}) /
\mathcal{D}_\ell^{\Lambda\rm CDM}$ together with the $1\sigma$ and
$2\sigma$ cosmic variance bands. The cosmic variance fractional
uncertainty is $\sigma / \mathcal{D}_\ell = \sqrt{2/[(2\ell+1)f_{\rm sky}]}$,
where $f_{\rm sky} \simeq 0.86$ for the Commander mask. At $\ell = 2$ this
gives $\sigma / \mathcal{D}_\ell \simeq 68\%$, decreasing to $\sim 20\%$
at $\ell = 30$.

Even the active-cutoff model, which produces the strongest suppression
in the CCR parameter space consistent with the data, remains within
the $1\sigma$ cosmic variance band at all multipoles $\ell \geq 3$.
This illustrates the fundamental limitation of the low-$\ell$ TT
spectrum for constraining infrared cutoff models: the irreducible
cosmic variance floor at low multipoles renders any suppression signal
statistically indistinguishable from a fluctuation of the standard model,
consistent with the near-unity Bayes factor reported in
section~\ref{sec:model_comparison}.

\begin{figure}[t]
  \centering
  \includegraphics[width=0.85\linewidth]{fig_cl_data_comparison}
  \caption{Low-$\ell$ TT power spectrum: Planck 2018 Commander data
    (grey points with error bars) compared with the $\Lambda$CDM prediction
    (black curve) and three CCR models. The grey band shows the
    $\pm 1\sigma$ cosmic variance envelope around $\Lambda$CDM. The
    CCR best-fit (green dashed) is indistinguishable from $\Lambda$CDM;
    the active-cutoff model (red) produces $\sim 50\%$ suppression at
    $\ell = 2$ but remains within the cosmic variance band for $\ell \geq 3$.}
  \label{fig:cl_data_comparison}
\end{figure}

\begin{figure}[t]
  \centering
  \includegraphics[width=0.85\linewidth]{fig_cl_residuals_cv}
  \caption{Fractional residuals of the CCR models relative to $\Lambda$CDM.
    The dark and light grey bands show the $1\sigma$ and $2\sigma$ cosmic
    variance, respectively. Planck Commander data residuals are shown as grey
    points. The cosmic variance at $\ell = 2$ is $\sim 68\%$, demonstrating
    why the low-$\ell$ power spectrum cannot distinguish the CCR suppression
    from a statistical fluctuation at current sensitivity.}
  \label{fig:cl_residuals}
\end{figure}

% ============================================================
% Section 4.6: Prior Sensitivity Analysis
% ============================================================

\subsection{Prior sensitivity analysis}
\label{subsec:prior_sensitivity}

To verify that our results are robust to the choice of prior distributions
on the CCR parameters, we perform an importance-reweighting analysis of the
MCMC chains under five alternative prior configurations, listed in
table~\ref{tab:prior_sensitivity}. These include narrowing and widening
the prior ranges on both $\log_{10}(\ln D)$ and $\alpha$, as well as a
joint restriction.

The reweighted posteriors confirm two key findings. First, the posterior
on $\log_{10}(\ln D)$ remains approximately uniform under all prior
choices, indicating that the data are uninformative with respect to the
Hilbert space dimension and the posterior simply tracks the prior.
Second, and more importantly, the standard $\Lambda$CDM parameters
are entirely unaffected by the choice of CCR priors: the shifts in
$H_0$, $n_s$, and $\tau$ are below $0.02\,\mathrm{km\,s^{-1}\,Mpc^{-1}}$,
$0.0002$, and $0.00002$, respectively, across all configurations. These
shifts are negligible compared to the posterior widths, confirming the
decoupling reported in section~\ref{subsec:decoupling}.

Figure~\ref{fig:prior_sensitivity} displays the reweighted posteriors
for $\log_{10}(\ln D)$, $\alpha$, and $H_0$, illustrating the stability
of the results.

\begin{table}[t]
  \centering
  \begin{tabular}{lccccc}
    \hline\hline
    Prior configuration & $\langle\log_{10}(\ln D)\rangle$
      & $\langle\alpha\rangle$ & $\langle H_0\rangle$
      & $\langle n_s\rangle$ & $N_{\rm eff}$ \\
    \hline
    Baseline: $[12, 16] \times [1, 6]$
      & $14.03 \pm 1.15$ & $3.44 \pm 1.44$ & $67.24 \pm 0.58$
      & $0.9647 \pm 0.0042$ & 5652 \\
    Narrow $\ln D$: $[13, 15] \times [1, 6]$
      & $14.01 \pm 0.57$ & $3.47 \pm 1.44$ & $67.23 \pm 0.59$
      & $0.9646 \pm 0.0042$ & 2862 \\
    Wide $\ln D$: $[11, 17] \times [1, 6]$
      & $14.03 \pm 1.15$ & $3.44 \pm 1.44$ & $67.24 \pm 0.58$
      & $0.9647 \pm 0.0042$ & 5652 \\
    Narrow $\alpha$: $[12, 16] \times [1, 4]$
      & $14.03 \pm 1.15$ & $2.47 \pm 0.85$ & $67.23 \pm 0.59$
      & $0.9647 \pm 0.0043$ & 3358 \\
    Both narrow: $[13, 15] \times [1, 4]$
      & $14.00 \pm 0.56$ & $2.50 \pm 0.85$ & $67.22 \pm 0.60$
      & $0.9646 \pm 0.0043$ & 1737 \\
    \hline\hline
  \end{tabular}
  \caption{Prior sensitivity analysis via importance reweighting. The
    mean and standard deviation of key parameters are shown for each
    prior configuration. The $\Lambda$CDM parameters ($H_0$, $n_s$) are
    stable to within $\Delta H_0 < 0.02$ and $\Delta n_s < 0.0002$
    across all choices, confirming robustness. $N_{\rm eff}$ is the
    effective sample size after reweighting.}
  \label{tab:prior_sensitivity}
\end{table}

\begin{figure}[t]
  \centering
  \includegraphics[width=\linewidth]{fig_prior_sensitivity}
  \caption{Prior sensitivity analysis. Reweighted posteriors for
    $\log_{10}(\ln D)$ (left), $\alpha$ (centre), and $H_0$ (right)
    under five alternative prior configurations. The $\log_{10}(\ln D)$
    posterior tracks the prior in all cases, confirming that the data are
    uninformative. The $H_0$ posterior is effectively identical across all
    prior choices, demonstrating complete decoupling of the CCR and
    $\Lambda$CDM parameter sectors.}
  \label{fig:prior_sensitivity}
\end{figure}
